% Section 15.8, from lectures/L15/appendix/b_exercises.md.

\section{Exercises}
\label{c15:sec:exercises}

\begin{aside}
These exercises consolidate the chapter. Design your own \file{rx_shift_reg.vhd} from the
specification in \secref{c15:sec:iface} to \ref{c15:sec:trace}; the testbench handed out, run per
appendix~\ref{app:sim}, is the check that it behaves correctly.
\end{aside}

\exercise{\code{rx_shift_reg}}{VHDL}
\label{c15:ex:rxsr}
Write your own \file{rx_shift_reg.vhd} from the chapter. Verify it:

\begin{shell}
cd controller
ghdl -a --std=93 can_def.vhd rx_shift_reg.vhd rx_shift_reg_tb.vhd
ghdl -e --std=93 rx_shift_reg_tb
ghdl -r --std=93 rx_shift_reg_tb --assert-level=error
\end{shell}

If a check fails, read the assertion message; it names exactly which sample it concerns and what was
expected.

\exercise{A deliberate stuffing violation}{Simulation}
\label{c15:ex:violation}
\xpart{a} Construct (on paper) a 6-bit sequence on the wire that constitutes a genuine violation of
the bit stuffing rule: five identical real bits in a row, followed by a sixth identical bit where a
stuff bit should have broken the run.

\xpart{b} Feed your sequence from \textbf{a)} into \code{rx_shift_reg} (either by extending
\file{rx_shift_reg_tb.vhd} or by writing a short new testbench) and confirm that \code{stuff_error}
pulses on exactly the sixth bit, not earlier.

\xpart{c} \code{rx_shift_reg} makes no attempt to recover or resynchronise after a
\code{stuff_error}; it merely reports it for one cycle and goes on accumulating as if nothing had
happened. Explain, in one or two sentences, why that is an acceptable design here, given how
\code{can_controller} (\chapref{c17:ch}) will react to it.

\exercise{The right-justification, and what sits above it}{Simulation}
\label{c15:ex:justify}
\Secref{c15:sec:iface} says that a completed group ends up \textbf{right-justified in its lowest
\code{bit_count} bits}, and that \code{can_controller} relies on that when it slices up reconstructed
fields. Tests 3 and 4 in the testbench handed out exercise it: a 6-bit group, then a 3-bit group, with
no reset in between, with \code{bit_count} changed between the groups in exactly the way
\code{can_controller} does between one field and the next.

Assume a just-cleared \code{rx_shift_reg}, \code{enable = '1'} and no stuff bits anywhere (no run ever
reaches five). The process in \secref{c15:sec:process} shifts every accepted bit into the register's
least significant end and never clears the register between groups; assume yours does the same. The
testbench handed out does not actually require that, and part \textbf{d)} returns to why.

\xpart{a} With \code{bit_count = 6}, feed in the six bits \code{0 1 0 1 0 1} (in that order, one per
\code{sample}). Write out all eight bits of \code{data} when \code{valid} pulses. What is in
\code{data(7 downto 6)}, and why?

\xpart{b} Set \code{bit_count = 3} without clearing, and feed in \code{0 1 1}. Write out all eight
bits of \code{data} when \code{valid} pulses again. What is now in \code{data(7 downto 3)}?

\xpart{c} Your answer to \textbf{b)} shows that the bits above \code{bit_count} hold residue from the
\emph{previous} group. Explain why that is harmless here, and what \code{can_controller} has to do
(and must never do) when it reads \code{rxsr_data} for a field narrower than 8 bits.

\xpart{d} \code{rx_shift_reg} holds two rather different kinds of state: the shift register
accumulating the current group, and the run counter (\code{consecutive}, \code{last_bit}) deciding
where stuff bits belong. Clearing them is not equally safe in both cases.
\begin{itemize}
  \item Suppose the module cleared its \textbf{shift register} at every \code{valid}. Which of your
    answers above would change, and would any field ever be \emph{misread} as a result? Note that the
    testbench handed out checks only the lowest \code{bit_count} bits of \code{data}, so both designs
    pass it.
  \item Suppose it cleared the \textbf{run counter} at the same moment. Name what breaks, and say
    which rule from \chapref{c14:ch} that mirrors.
\end{itemize}

The difference between the two answers is why \secref{c15:sub:defaults} says the run counter is
cleared only by reset and says nothing at all about the shift register.
